\documentclass[12pt,a4paper]{article}

\usepackage{fontspec}
\usepackage[french]{babel}
\setmainfont{EBGaramond}
\usepackage[a4paper,margin=24mm]{geometry}
\usepackage{microtype}

\usepackage{csquotes}
% \textelp{}
% https://tex.stackexchange.com/a/294634/295527

\begin{document}\thispagestyle{empty}

\subsection*{De la grandeur de la Terre.}

37. L'observation de la hauteur du pôle et de la hauteur de l'équateur, ou, si l'on veut, de la hauteur méridienne du soleil en différents pays, fut la première chose qui dut apprendre aux hommes que la terre était ronde. Ce fut d'abord par l'ombre du soleil que l'on détermina les différences de hauteurs du pôle ; plus on avançait vers le nord, plus ces ombres mesurées le même jour se trouvaient longues ; ce qui prouvait que la hauteur du soleil au-dessus de l'horizon était devenue plus petite, et que l'observateur situé vers le nord n'était pas sur le même plan que l'observateur situé vers le midi. On dut en conclure que la terre était arrondie.

38. L'ombre de la terre dans les éclipses de lune paraît toujours ronde ; les vaisseaux vus de loin en pleine mer, disparaissent par degrés ; et on les voit descendre et se perdre peu à peu, par la courbure de la surface des eaux. Telles furent les marques auxquelles les anciens philosophes reconnurent la courbure et la rondeur de la terre.

39. Après avoir ainsi reconnu la rondeur de la terre, on se servit du même moyen pour connaître sa grandeur : et le changement des latitudes et des hauteurs, soit du pôle, soit des astres, servit à connaître l'étendue de notre globe en en mesurant une petite partie. \textelp{}

40. Autre exemple : on trouve en allant vers le nord que la latitude d'Amiens est plus grande que celle de Paris d'un degré, ou que le soleil à midi est d'un degré plus bas à Amiens qu'à Paris ; c'est une preuve que la terre a un degré de courbure depuis Paris jusqu'à Amiens ; or cette distance mesurée en allant toujours du midi au nord, s'est trouvée de 25 lieues, chaune de 2283 toises ; donc un degré de la terre, ou la 360e partie de toute sa circonférence, a 25 lieues d'étendue ; d'où il suit que la circonférence entière ou le tour de la terre vaut 9000 lieues ; car 25 fois 360 font 9000.

\begin{flushright}\large
\textsc{lalande}, \textit{Abrégé d'astronomie}, Livre premier.
\end{flushright}
\end{document}
